%% toc: P2 11 11.4,11.5,11.6
\section{Proofs for duality for arbitrary propositional formulas}
\label{sec:annex-p02-w02-p02-010-duality-arbitrary}

Throughout, the complemented valuation \(\bar v\), the complementing
substitution \(\upsilon_{\neg}\) and NNF duals are as in
Convention~\ref{def:p02-w02-p02-010-complement}.  The negation normal form
theorem is the unlabelled theorem of that name in
Section~\ref{sec:prop-syntax-semantics-04}.  The NNF formulas form the least
set of formulas containing every literal, \(\bottom\) and \(\topv\) and closed
under forming \((\beta\land\gamma)\) and \((\beta\lor\gamma)\) from members
\(\beta,\gamma\); this is the definition of negation normal form in
Section~\ref{sec:prop-syntax-semantics-04}.

\begin{annexitem}{annex:p02-w02-p02-010-nnf-presentation}{thm:p02-w02-p02-010-nnf-presentation}{Unique NNF presentation}
\begin{proof}
Write \(\mathsf{NNF}\) for the set of NNF formulas, and let \(\mathsf{NNF}'\) be
the set of members of \(\mathsf{NNF}\) that are a variable, a negated variable,
\(\bottom\), \(\neg\bottom\), or of the form \((\beta\land\gamma)\) or
\((\beta\lor\gamma)\) with \(\beta,\gamma\in\mathsf{NNF}\).  Every literal, and
\(\bottom\) and \(\topv=\neg\bottom\), lies in \(\mathsf{NNF}'\).  If
\(\beta,\gamma\in\mathsf{NNF}\), then \((\beta\land\gamma)\) and
\((\beta\lor\gamma)\) lie in \(\mathsf{NNF}\) and have the shape required for
\(\mathsf{NNF}'\).  So \(\mathsf{NNF}'\) contains the generators of \(\mathsf{NNF}\)
and is closed under its two constructions; since \(\mathsf{NNF}\) is the least
such set, \(\mathsf{NNF}\subseteq\mathsf{NNF}'\).  Hence every NNF formula has one
of the six listed shapes, and in the two binary shapes its components are NNF
formulas.

The six shapes are mutually exclusive by
Theorem~\ref{thm:prop-unique-readability}.  A variable, \(\bottom\), a
negation and a binary formula satisfy four different formation clauses.  The
two negation shapes \(\neg p_n\) and \(\neg\bottom\) have different immediate
subformulas, \(p_n\ne\bottom\), and the immediate subformula of a negation is
unique.  The main connective of a binary formula is unique, so no formula is
both a conjunction and a disjunction, and its two immediate subformulas are
unique; these are the \(\beta\) and \(\gamma\) of the lemma.

Therefore each NNF formula \(\alpha\) falls under exactly one clause of
Definition~\ref{def:p02-nnf-dual}, with uniquely determined components, which
are NNF formulas strictly shorter than \(\alpha\).  Recursion on the length of
\(\alpha\) therefore defines \(\alpha^{\mathrm d}\) uniquely.  Induction on the
length shows that \(\alpha^{\mathrm d}\) is an NNF formula: a literal is sent
to itself, \(\bottom\) to \(\topv\) and \(\topv\) to \(\bottom\), all NNF
formulas; and \((\beta\land\gamma)^{\mathrm d}=(\beta^{\mathrm d}\lor\gamma^{\mathrm d})\)
and \((\beta\lor\gamma)^{\mathrm d}=(\beta^{\mathrm d}\land\gamma^{\mathrm d})\)
are NNF formulas because \(\beta^{\mathrm d}\) and \(\gamma^{\mathrm d}\) are,
by the induction hypothesis, and \(\mathsf{NNF}\) is closed under both binary
constructions.
\end{proof}
\end{annexitem}

\begin{annexitem}{annex:p02-w02-p02-010-arbitrary-duality}{thm:p02-w02-p02-010-arbitrary-duality}{Duality for arbitrary formulas}
\begin{proof}
\emph{Clause~1.}  By the negation normal form theorem there is an NNF formula
\(\alpha\) with \(\alpha\equiv\varphi\).  By
Lemma~\ref{thm:p02-w02-p02-010-nnf-presentation}, \(\alpha^{\mathrm d}\) is
defined, so it is an NNF dual of \(\varphi\).

\emph{Clause~2.}  Let \(\alpha\) be an NNF formula with \(\alpha\equiv\varphi\)
and let \(v\) be a valuation.  By
Lemma~\ref{thm:p02-w02-p02-010-nnf-presentation}, \(\alpha\) has exactly one
NNF presentation, and \(\alpha^{\mathrm d}\) is its dual.  The semantic
identity of Theorem~\ref{thm:p02-propositional-duality}, proved in part~1 of
Annex~\ref{annex:p2.1.8}, gives \(v\entails\alpha^{\mathrm d}\) iff
\(\bar v\not\entails\alpha\).  Because \(\alpha\equiv\varphi\), the formulas
\(\alpha\) and \(\varphi\) take the same value under \(\bar v\); hence
\(v\entails\alpha^{\mathrm d}\) iff \(\bar v\not\entails\varphi\).

For the second form, Theorem~\ref{thm:p02-uniform-substitution} gives
\(\widehat v(\varphi[\upsilon_{\neg}])=\widehat{v_{\upsilon_{\neg}}}(\varphi)\),
and \(v_{\upsilon_{\neg}}=\bar v\) by
Convention~\ref{def:p02-w02-p02-010-complement}.  So, by the negation clause,
\(v\entails\neg(\varphi[\upsilon_{\neg}])\) iff
\(\widehat v(\varphi[\upsilon_{\neg}])=\mathbf F\) iff
\(\widehat{\bar v}(\varphi)=\mathbf F\) iff \(\bar v\not\entails\varphi\) iff,
by the first form, \(v\entails\alpha^{\mathrm d}\).  This holds for every
\(v\), so \(\alpha^{\mathrm d}\equiv\neg(\varphi[\upsilon_{\neg}])\).
Conversely, reading the same chain of equivalences from the right, the second
form gives back \(v\entails\alpha^{\mathrm d}\) iff
\(\bar v\not\entails\varphi\) for every \(v\).

\emph{Clause~3.}  If \(\alpha\) and \(\alpha'\) are NNF formulas logically
equivalent to \(\varphi\), clause~2 gives
\(\alpha^{\mathrm d}\equiv\neg(\varphi[\upsilon_{\neg}])\) and
\(\alpha'^{\mathrm d}\equiv\neg(\varphi[\upsilon_{\neg}])\).  Logical
equivalence is symmetric and transitive, since it means having the same value
under every valuation, so \(\alpha^{\mathrm d}\equiv\alpha'^{\mathrm d}\).  If
\(\varphi\) is an NNF formula, then \(\varphi\equiv\varphi\) and
\(\varphi^{\mathrm d}\) is an NNF dual of \(\varphi\) by definition.

\emph{Clause~4.}  Let \(\alpha=\alpha_0^{\mathrm d}\) and
\(\beta=\beta_0^{\mathrm d}\) with \(\alpha_0,\beta_0\) NNF formulas,
\(\alpha_0\equiv\varphi\) and \(\beta_0\equiv\psi\).  By clause~2, for every
valuation \(v\): \(v\entails\alpha\) iff \(\bar v\not\entails\varphi\), and
\(v\entails\beta\) iff \(\bar v\not\entails\psi\).

Suppose \(\varphi\entails\psi\) and let \(v\entails\beta\).  Then
\(\bar v\not\entails\psi\).  If \(\bar v\) satisfied \(\varphi\), it would
satisfy \(\psi\); so \(\bar v\not\entails\varphi\), and \(v\entails\alpha\).
Hence \(\beta\entails\alpha\).

Suppose \(\beta\entails\alpha\) and let \(u\entails\varphi\).  Put
\(v=\bar u\), so that \(\bar v=u\).  Then \(\bar v\entails\varphi\), so
\(v\not\entails\alpha\).  If \(v\) satisfied \(\beta\), it would satisfy
\(\alpha\); so \(v\not\entails\beta\), hence \(\bar v\entails\psi\), that is,
\(u\entails\psi\).  Hence \(\varphi\entails\psi\).

By the definitions of Section~\ref{sec:prop-syntax-semantics-02},
\(\varphi\equiv\psi\) holds iff every valuation satisfying either formula
satisfies the other, that is, iff \(\varphi\entails\psi\) and
\(\psi\entails\varphi\).  By the two implications just proved, applied in both
directions, this holds iff \(\beta\entails\alpha\) and
\(\alpha\entails\beta\), that is, iff \(\alpha\equiv\beta\).

\emph{Clause~5.}  Let \(\alpha\) be an NNF dual of \(\varphi\).  By clause~2,
\(v\entails\alpha\) iff \(\bar v\not\entails\varphi\), for every \(v\).  As
\(v\) ranges over all valuations, so does \(\bar v\), since \(u=\bar v\) for
\(v=\bar u\).  Hence \(\varphi\) is valid iff \(\bar v\entails\varphi\) for
every \(v\) iff \(v\not\entails\alpha\) for every \(v\) iff \(\alpha\) is
unsatisfiable.  Likewise \(\varphi\) is unsatisfiable iff
\(\bar v\not\entails\varphi\) for every \(v\) iff \(v\entails\alpha\) for every
\(v\) iff \(\alpha\) is valid.
\end{proof}
\end{annexitem}

\begin{annexitem}{annex:p02-w02-p02-010-normal-forms}{thm:p02-w02-p02-010-normal-forms}{Normal forms of the dual}
\begin{proof}
\emph{Clause~1.}  A formula in conjunctive normal form is an NNF formula: it
is \(\topv\), or a right-associated conjunction of formulas
\(\operatorname{form}_{\mathsf{CPL}}(C)\), each of which is \(\bottom\) or a
right-associated disjunction of literals.  So if \(\gamma\) is in conjunctive
normal form and \(\gamma\equiv\varphi\), then \(\gamma^{\mathrm d}\) is an NNF
dual of \(\varphi\), and part~6 of Annex~\ref{annex:p2.1.8} shows that
\(\gamma^{\mathrm d}\) is in disjunctive normal form.  Such a \(\gamma\)
exists by Theorem~\ref{thm:prop-normal-forms}.

\emph{Clause~2.}  A formula in disjunctive normal form is an NNF formula: it is
a right-associated disjunction (\(\bottom\) when empty) of right-associated
conjunctions of literals (\(\topv\) when empty).  So if \(\delta\) is in
disjunctive normal form and \(\delta\equiv\varphi\), then
\(\delta^{\mathrm d}\) is an NNF dual of \(\varphi\), and part~7 of
Annex~\ref{annex:p2.1.8} shows that \(\delta^{\mathrm d}\) is logically
equivalent to the conjunctive normal form built from the clauses
\(\{L_{i1},\ldots,L_{ik_i}\}\); when \(n=0\), \(\delta=\bottom\) and
\(\delta^{\mathrm d}=\topv\), which is that conjunctive normal form.  Part~8 of
Annex~\ref{annex:p2.1.8} proves that \((p_1\lor p_1)\), the dual of the
disjunctive normal form \((p_1\land p_1)\), is not in conjunctive normal form.
\end{proof}
\end{annexitem}
